%% Question bank chapter file - generated by bank_tools.py build (do not edit while the bank is still being built:
%% it is regenerated from the per-test data in Question Bank\data). Every question is one bankq block:
%%   \begin{bankq}[id=..., source={...}, paper=.., qno=.., type=..., subtopic={...}, answer={...}, answer_src=key|solved,
%%                 working={...}, difficulty=1-10, diff_parts={K S I C F}, diff_note={...},
%%                 group=<passage id or empty>, also_in={duplicate copies in other tests}]
%% difficulty = min(10, K+S+I+C+F-1): K concepts 1-3, S steps 1-4, I insight 0-2, C computation 0-2, F format 0-1
%%   ...question LaTeX, options, figures (\bfig / \bfigi -> img/)...
%%   \end{bankq}
%% Passages shared by several questions are bankpassage blocks; their questions carry group=<passage id>.
\documentclass[10pt]{article}
\usepackage[a4paper,margin=1.8cm]{geometry}
\usepackage[T1]{fontenc}
\usepackage{amsmath,amssymb,mathtools}
\usepackage{graphicx}
\usepackage[export]{adjustbox}
\usepackage{xcolor,needspace,keyval}
\usepackage{enumitem}
\usepackage{booktabs,array,tabularx}
\usepackage[hidelinks]{hyperref}
\setlength{\parindent}{0pt}
\setlength{\parskip}{3pt}
\makeatletter
\define@key{bq}{id}{\def\bq@id{#1}}\define@key{bq}{source}{\def\bq@src{#1}}\define@key{bq}{paper}{\def\bq@paper{#1}}
\define@key{bq}{qno}{\def\bq@qno{#1}}\define@key{bq}{type}{\def\bq@type{#1}}\define@key{bq}{subtopic}{\def\bq@sub{#1}}
\define@key{bq}{answer}{\def\bq@ans{#1}}\define@key{bq}{answer_src}{\def\bq@asrc{#1}}\define@key{bq}{working}{\def\bq@work{#1}}
\define@key{bq}{group}{\def\bq@group{#1}}\define@key{bq}{also_in}{\def\bq@also{#1}}
\define@key{bq}{difficulty}{\def\bq@diff{#1}}\define@key{bq}{diff_parts}{\def\bq@dparts{#1}}\define@key{bq}{diff_note}{\def\bq@dnote{#1}}
\newenvironment{bankq}[1][]{\def\bq@work{}\def\bq@also{}\def\bq@group{}\def\bq@diff{?}\def\bq@dparts{}\def\bq@dnote{}\setkeys{bq}{#1}%
  \par\needspace{4\baselineskip}\medskip\noindent{\small\textbf{\bq@id}\hfill\color{gray}[\bq@src, Paper \bq@paper, Q\bq@qno]
  \bq@type\ $\cdot$ \bq@sub\ $\cdot$ \textbf{D\bq@diff}}\par\nopagebreak\relax}
  {\par{\small\color{teal}Answer: \textbf{\bq@ans}\ifx\bq@work\@empty\else\ (solved: \bq@work)\fi
  \ifx\bq@also\@empty\else\par\color{gray}Also in: \bq@also\fi}\par}
\define@key{bp}{id}{\def\bp@id{#1}}\define@key{bp}{source}{\def\bp@src{#1}}\define@key{bp}{questions}{\def\bp@qs{#1}}
\newenvironment{bankpassage}[1][]{\setkeys{bp}{#1}\par\needspace{4\baselineskip}\bigskip\noindent
  {\small\textbf{Passage \bp@id}\hfill\color{gray}[\bp@src; questions \bp@qs]}\par\nopagebreak\relax\itshape}{\par}
\makeatother
\newcommand{\bfig}[2]{\par\smallskip{\centering\includegraphics[width=#1pt,max width=0.62\linewidth,max height=0.32\textheight]{img/#2}\par}\smallskip}
\newcommand{\bfigi}[2]{\includegraphics[width=#1pt,max width=\linewidth,max height=0.32\textheight,valign=c]{img/#2}}
\newcommand{\bmissing}[1]{\textit{[figure missing in the source conversion: \detokenize{#1}]}}
\newenvironment{qoptions}{\begin{itemize}[leftmargin=1.8em,itemsep=1pt,topsep=2pt,label={}]}{\end{itemize}}
\newcommand{\opt}[2]{\item[(#1)] #2}
\providecommand{\fwstart}{}\providecommand{\fwend}{}\providecommand{\qfullsplit}{}
\newcolumntype{L}[1]{>{\raggedright\arraybackslash}p{#1}}
\newcolumntype{C}[1]{>{\centering\arraybackslash}p{#1}}
\newcolumntype{R}[1]{>{\raggedleft\arraybackslash}p{#1}}
\providecommand{\sectionhdr}[1]{}\providecommand{\subsectionhdr}[1]{}


\begin{document}
{\Large\bfseries Logarithms --- Question Bank}\par\smallskip
{\small Mathematics chapter M10. Questions from all converted test papers, grouped by source series. Answers: printed key where it matches the paper, otherwise solved (answer\_src=solved, with working). Exact repeats across papers appear once, with also\_in listing the other copies.}\par\medskip\hrule

\section*{TOAS papers}
\begin{bankq}[id=TOAS01-P1-MATH-S2-Q09, source={TOAS - Test 01 (30 Jan 2022)}, paper=1, qno=9, type=multi, subtopic={Properties}, answer={A,B}, answer_src=key, difficulty=2, diff_parts={K1 S1 I0 C0 F1}, diff_note={concepts: properties of logarithms | steps: use a = b - 2, c = b + 2 in each option}]
If $a,b,c$ are three consecutive even natural number, then which of the following is (are) true -
\begin{qoptions}
\opt{A}{$\log_{10}(4+ac)=2\log_{10}b$}
\opt{B}{$\log_{10}(a+c)=\log_{10}2+\log_{10}b$}
\opt{C}{$\log_{10}(a^2+c^2)=3\log_{10}2+\log_{10}b$}
\opt{D}{$\log_{10}(a+b+c)=\log_{10}6+\log_{10}b$}
\end{qoptions}
\end{bankq}
\begin{bankq}[id=TOAS01-P2-MATH-S2-Q10, source={TOAS - Test 01 (30 Jan 2022)}, paper=2, qno=10, type=numerical, subtopic={Equations and inequalities}, answer={7}, answer_src=solved, working={Need $(a-5)x^2+4x+(a-5)\ge0$ for all $x$: $a>5$ and $16\le4(a-5)^2$, $a\ge7$; smallest integer 7}, difficulty=4, diff_parts={K2 S2 I0 C0 F1}, diff_note={concepts: logarithmic inequality, quadratic nonnegativity | steps: convert to (a - 5)x\^{}2 + 4x + (a - 5) >= 0; leading coefficient positive; discriminant condition; smallest integer}]
The smallest integral value of $a$ for which the inequality $1+\log_5(1+x^2)\le\log_5(ax^2+4x+a)$ is true for all $x\in R$ is
\end{bankq}
\begin{bankq}[id=TOAS02-P1-MATH-S1-Q06, source={TOAS - Test 02 (13 Feb 2022)}, paper=1, qno=6, type=single, subtopic={Equations and inequalities}, answer={C}, answer_src=key, difficulty=4, diff_parts={K2 S2 I1 C0 F0}, diff_note={concepts: logarithms, modulus equation | steps: express both logarithms in log\_2 x; relate the three absolute value terms; use the equality case of the triangle inequality | I: equality condition of the triangle inequality}]
Solution set of equation is $\left|1-\log_{\frac16}x\right|+|\log_2x|+2=\left|3-\log_{\frac16}x+\log_{\frac12}x\right|$, $a,b\in N$, then the value of $(a+b)$ is
\begin{qoptions}
\opt{A}{5}
\opt{B}{6}
\opt{C}{7}
\opt{D}{8}
\end{qoptions}
\end{bankq}
\begin{bankq}[id=TOAS02-P1-MATH-S2-Q10, source={TOAS - Test 02 (13 Feb 2022)}, paper=1, qno=10, type=multi, subtopic={Equations and inequalities}, answer={A,C,D}, answer_src=key, difficulty=6, diff_parts={K2 S2 I2 C0 F1}, diff_note={concepts: AM-GM inequality, completing the square | steps: apply AM-GM to the two exponentials; complete the squares in the exponent; equality needs x = y = -1/2; check the options | I: AM-GM equality}]
Let ordered pair $(x,y)$ satisfy the equation $49^{x^2+y}+49^{y^2+x}=\dfrac{2}{\sqrt7}$ then which of the following is/are correct ?
\begin{qoptions}
\opt{A}{Only 1 ordered pair of $(x,y)$ is possible}
\opt{B}{More than 1 ordered pair of $(x,y)$ are possible}
\opt{C}{In one of the ordered pair $|x|+|y|=1$}
\opt{D}{In one of the ordered pair $x-y=0$}
\end{qoptions}
\end{bankq}
\begin{bankq}[id=TOAS02-P1-MATH-S3-Q03, source={TOAS - Test 02 (13 Feb 2022)}, paper=1, qno=3, type=numerical, subtopic={Properties}, answer={1}, answer_src=solved, working={With $A=\ln a$ etc.: $\frac{A^3+B^3+C^3}{ABC}=3$, so $A+B+C=0$ and $abc=1$}, difficulty=6, diff_parts={K2 S3 I1 C0 F1}, diff_note={concepts: logarithm base change, symmetric identity | steps: convert to natural logarithms; simplify each product; combine into (A\^{}3 + B\^{}3 + C\^{}3)/(A B C) = 3; use the factorisation; A + B + C = 0 | I: a\^{}3 + b\^{}3 + c\^{}3 - 3abc}]
Let $a,b$ and $c$ are distinct positive numbers, none of them is equal to unity such that $\log_ba.\log_ca+\log_ab.\log_cb+\log_ac.\log_bc-\log_ba\sqrt a\log_{\sqrt c}b^{1/3}\log_ac^3=0$, then the value of $abc$ is -
\end{bankq}
\begin{bankq}[id=TOAS03-P1-MATH-S2-Q12, source={TOAS - Test 03 (27 Feb 2022)}, paper=1, qno=12, type=multi, subtopic={Properties}, answer={A,B,C,D}, answer_src=key, difficulty=2, diff_parts={K1 S1 I0 C0 F1}, diff_note={concepts: signs of logarithms | steps: each of the four expressions is positive}]
Which of the following is (are) positive -
\begin{qoptions}
\opt{A}{$\log_7 25$}
\opt{B}{$\log_3(\log_3 9)$}
\opt{C}{$\dfrac{\log_3 5}{\log_3(2+x^2)}$}
\opt{D}{$\left(\log_{1/2}3\right)^{10}$}
\end{qoptions}
\end{bankq}
\begin{bankq}[id=TOAS19-P1-MATH-S3-Q05, source={TOAS - Test 19 (05 Feb 2023)}, paper=1, qno=5, type=numerical, subtopic={Equations and inequalities}, answer={5}, answer_src=solved, working={With $t=\log_{10}x$: $(t/3)^{t^2-3t+3}=t/3$. Base 1 ($t=3$); exponent 1 ($t=1,2$); base $-1$ with odd exponent 21 ($t=-3$); base 0 with positive exponent ($t=0$): 5 solutions}, difficulty=6, diff_parts={K1 S2 I2 C1 F1}, diff_note={concepts: exponential equations with logarithmic bases | steps: substitute t; cases for base 1 or exponent 1; negative base case; count}, also_in={TNAS - Test 01 (05 Feb 2023) P2 Q5}]
The number of real solutions of the equation $\left(\dfrac{\log_{10}x}3\right)^{\log_{10}^2x-3\log_{10}x+3}=\log_{10}\sqrt[3]x$, is
\end{bankq}

\section*{TOPS papers}
\begin{bankq}[id=TOPS01-P1-MATH-S1-Q06, source={TOPS - Test 01 (31 Jul 2022)}, paper=1, qno=6, type=multi, subtopic={Properties}, answer={A,B,C}, answer_src=key, difficulty=6, diff_parts={K2 S2 I1 C1 F1}, diff_note={concepts: logarithm properties, inequalities | steps: change of base on each; bound the product; compare the bases | four log comparisons}]
Which of the following hold(s) good ?
\begin{qoptions}
\opt{A}{$(\log_35)(\log_34)(\log_{10}11)>1$}
\opt{B}{$(\log_23)(\log_{1/3}5)(\log_46)<1$}
\opt{C}{$\log_{100}7>\log_{29}3$}
\opt{D}{$\log_\pi5>\log_e5$}
\end{qoptions}
\end{bankq}
\begin{bankq}[id=TOPS01-P1-MATH-S3-Q03, source={TOPS - Test 01 (31 Jul 2022)}, paper=1, qno=3, type=numerical, subtopic={Properties}, answer={-4.00}, answer_src=key, difficulty=6, diff_parts={K2 S2 I1 C1 F1}, diff_note={concepts: logarithm properties, perfect square | steps: let t equal log y over log x; write the terms; recognise the square; simplify}]
The value of
\par\noindent\hfill$\begin{aligned}\sqrt{(\log_xy)\left(\log_x\left(x^2y\right)\right)\left(\log_x\left(x^4y\right)\right)\left(\log_x\left(x^6y\right)\right)+16}\\&-\left(\log_x\left(x^2y\right)\right)\left(\log_x\left(x^4y\right)\right)\end{aligned}$\hfill\null\par
$x>1$, $y>1$; is
\end{bankq}
\begin{bankq}[id=TOPS02-P1-MATH-S1-Q03, source={TOPS - Test 02 (28 Aug 2022)}, paper=1, qno=3, type=single, subtopic={Properties}, answer={A}, answer_src=key, difficulty=4, diff_parts={K1 S2 I1 C1 F0}, diff_note={concepts: logarithm properties | steps: express the numbers in terms of a and b; simplify; match}]
If $\log_{10}25=a$, $\log_{10}225=b$ then the value of $\log_{10}\left(\dfrac19\right)^2+\log_{10}\left(\dfrac1{2250}\right)$ is
\begin{qoptions}
\opt{A}{$2a-3b-1$}
\opt{B}{$2a+3b$}
\opt{C}{$2a+3b+1$}
\opt{D}{$2a+3b-1$}
\end{qoptions}
\end{bankq}
\begin{bankq}[id=TOPS02-P1-MATH-S2-Q07, source={TOPS - Test 02 (28 Aug 2022)}, paper=1, qno=7, type=multi, subtopic={Properties}, answer={A,C,D}, answer_src=key, difficulty=6, diff_parts={K2 S2 I1 C1 F1}, diff_note={concepts: logarithm properties, inequalities | steps: sign of ln sin x on the interval; compare the bases; check each option}]
Which of the following is/are true ?
\begin{qoptions}
\opt{A}{$(\tan x)^{\ell n(\sin x)}>(\cot x)^{\ell n(\sin x)}\ \forall x\in\left(0,\dfrac\pi4\right)$}
\opt{B}{$4^{\ell n(\operatorname{cosec}x)}>5^{\ell n(\operatorname{cosec}x)}\ \forall x\in\left(0,\dfrac\pi2\right)$}
\opt{C}{$\left(\dfrac12\right)^{\ell n(\cos x)}<\left(\dfrac13\right)^{\ell n(\cos x)}$, $\forall x\in\left(0,\dfrac\pi2\right)$}
\opt{D}{$2^{\ell n(\tan x)}>2^{\ell n(\sin x)}$, $\forall x\in\left(0,\dfrac\pi2\right)$}
\end{qoptions}
\end{bankq}
\begin{bankq}[id=TOPS02-P1-MATH-S2-Q12, source={TOPS - Test 02 (28 Aug 2022)}, paper=1, qno=12, type=multi, subtopic={Equations and inequalities}, answer={A,D}, answer_src=key, difficulty=8, diff_parts={K2 S3 I2 C1 F1}, diff_note={concepts: logarithm equation, quadratics | steps: simplify the exponent terms; reduce to a quadratic; check the domain; sum of solutions; product of solutions | log identity}]
If $S$ is sum of solutions \& $P$ is product of solutions of the equation $2^{\log_2(x^2-x)}-4x+5^{\log_{10}6}.2^{\log_{10}6}=0$, then -
\begin{qoptions}
\opt{A}{$S+P=11$}
\opt{B}{$S^2=16$}
\opt{C}{$P^2=25$}
\opt{D}{$S^2=25$}
\end{qoptions}
\end{bankq}
\begin{bankq}[id=TOPS02-P1-MATH-S3-Q06, source={TOPS - Test 02 (28 Aug 2022)}, paper=1, qno=6, type=numerical, subtopic={Equations and inequalities}, answer={7.00}, answer_src=key, difficulty=6, diff_parts={K2 S2 I1 C1 F1}, diff_note={concepts: logarithm properties, sum of squares | steps: both squares must vanish; find x; find y; evaluate | equation as sum of squares}]
If $\log_7^2(x+7)+\log_7^2(3+\log_2y)=0$, then value of $|x|+4y$ is
\end{bankq}
\begin{bankq}[id=TOPS02-P2-MATH-S1-Q06, source={TOPS - Test 02 (28 Aug 2022)}, paper=2, qno=6, type=multi, subtopic={Equations and inequalities}, answer={C,D}, answer_src=key, difficulty=9, diff_parts={K2 S3 I2 C2 F1}, diff_note={concepts: logarithm equation, domain | steps: drop the outer logarithm; drop the inner one; solve the quadratic; check bases; check arguments | nested logarithms}]
For the equation
\par\noindent\hfill$\log_{x^2+6x+8}\left(\log_{2x^2+2x+3}\left(x^2-2x\right)\right)=0\ -$\hfill\null\par
\begin{qoptions}
\opt{A}{number of solutions is two}
\opt{B}{sum of solutions is `$-4$'}
\opt{C}{number of solution is one}
\opt{D}{sum of solutions is `$-1$'}
\end{qoptions}
\end{bankq}
\begin{bankq}[id=TOPS04-P2-MATH-S1-Q19, source={TOPS - Test 04 (18 Dec 2022)}, paper=2, qno=19, type=single, subtopic={Equations and inequalities}, answer={B}, answer_src=key, difficulty=2, diff_parts={K1 S1 I0 C1 F0}, diff_note={concepts: logarithm equation | steps: combine the logarithms; solve}]
If $\log_3x^4=\log_3x^3+\log_37-\log_32$, then $x$ is equal to
\begin{qoptions}
\opt{A}{3.25}
\opt{B}{3.50}
\opt{C}{4.75}
\opt{D}{6.25}
\end{qoptions}
\end{bankq}
\begin{bankq}[id=TOPS05-P2-MATH-S2-Q10, source={TOPS - Test 05 (14 Jan 2023)}, paper=2, qno=10, type=multi, subtopic={Properties}, answer={A,B,C}, answer_src=key, difficulty=7, diff_parts={K2 S2 I2 C1 F1}, diff_note={concepts: logarithms, powers | steps: express a, b, c in one variable; solve the sum; values of logs}]
If $\log_ba=2$, $\log_cb=3$ and $a+b+c=27+4\sqrt3$ then value of :
\begin{qoptions}
\opt{A}{$\log_3(abc)=5$}
\opt{B}{$\log_3(a^2bc)=8$}
\opt{C}{$\log_3(abc^3)=6$}
\opt{D}{$\log_3(a^2bc^2)=6$}
\end{qoptions}
\end{bankq}

\section*{TEAS papers}
\begin{bankq}[id=TEAS01-P1-MATH-S1-Q06, source={TEAS - Test 01 (16 Apr 2023)}, paper=1, qno=6, type=multi, subtopic={Properties}, answer={A,B,C}, answer_src=key, difficulty=6, diff_parts={K2 S2 I1 C1 F1}, diff_note={concepts: logarithms, digits | steps: digits by log; change of base; solve for log of 2 base 3; check}]
Which of the following is/are correct ($\log_{10}2=0.3010$, $\log_{10}3=0.4771$)
\begin{qoptions}
\opt{A}{Number of digits in $(243)^{10}$ is 24}
\opt{B}{Number of digits in $2^{100}$ is 31}
\opt{C}{If $\log_45=a$ and $\log_56=b$ then $\log_32$ is $\dfrac1{2ab-1}$}
\opt{D}{If $\log_45=a$ and $\log_56=b$ then $\log_32=2ab-1$}
\end{qoptions}
\end{bankq}
\begin{bankq}[id=TEAS01-P1-MATH-S3-Q02, source={TEAS - Test 01 (16 Apr 2023)}, paper=1, qno=2, type=numerical, subtopic={Equations and inequalities}, answer={0.25}, answer_src=key, difficulty=7, diff_parts={K2 S2 I2 C1 F1}, diff_note={concepts: logarithm, roots | steps: substitute t; sum and product of roots; expression in the roots; value}]
If $\alpha$ and $\beta$ are roots of equation $(\log_2x)^2-3\log_2x+4=0$, then the value of $\log_\alpha\beta+\log_\beta\alpha$ is equal to -
\end{bankq}

\section*{TNAS papers}
\begin{bankq}[id=TNAS01-P1-MATH-S1-Q07, source={TNAS - Test 01 (05 Feb 2023)}, paper=1, qno=7, type=single, subtopic={Equations and inequalities}, answer={B}, answer_src=key, difficulty=5, diff_parts={K1 S2 I2 C1 F0}, diff_note={concepts: logarithmic inequality with trigonometric base | steps: base conditions; inequality; solve}]
The solution set of equation $\log_{\cos x}\left(\dfrac{2\sin x+3}{4}\right)\geqslant2$, where $x\in[0,\pi]$, is
\begin{qoptions}
\opt{A}{$\left[0,\dfrac\pi5\right]$}
\opt{B}{$\left(0,\dfrac{\pi}{10}\right]$}
\opt{C}{$\left[0,\dfrac{\pi}{10}\right]$}
\opt{D}{$\phi$ (null set)}
\end{qoptions}
\end{bankq}
\begin{bankq}[id=TNAS01-P1-MATH-S1-Q19, source={TNAS - Test 01 (05 Feb 2023)}, paper=1, qno=19, type=single, subtopic={Equations and inequalities}, answer={C}, answer_src=key, difficulty=4, diff_parts={K1 S2 I1 C1 F0}, diff_note={concepts: logarithm system | steps: change of base; solve for logs; product}]
If $\log_{27}x+\log_9y=\dfrac72$, $\log_{27}y+\log_9x=\dfrac23$, then the value of $xy$ is
\begin{qoptions}
\opt{A}{$27$}
\opt{B}{$81$}
\opt{C}{$243$}
\opt{D}{$216$}
\end{qoptions}
\end{bankq}
\begin{bankq}[id=TNAS07-P1-MATH-S1-Q17, source={TNAS - Test 07 (20 Aug 2023)}, paper=1, qno=17, type=single, subtopic={Equations and inequalities}, answer={A}, answer_src=key, difficulty=7, diff_parts={K1 S3 I2 C2 F0}, diff_note={concepts: number of solutions of a nested modulus equation | steps: case on the sign of x; graph of the log; graph of the modulus; count intersections; symmetry}]
The number of solutions of equation $\left|\ell n|x||-||x|-1|\right|=0$ is
\begin{qoptions}
\opt{A}{2}
\opt{B}{4}
\opt{C}{8}
\opt{D}{10}
\end{qoptions}
\end{bankq}

\section*{Star Review (Other) papers}
\begin{bankq}[id=STAR01-P2-MATH-S1-Q13, source={Star Review (Other) - Test 01 (12 Sep 2022)}, paper=2, qno=13, type=single, subtopic={Equations and inequalities}, answer={C}, answer_src=key, difficulty=5, diff_parts={K1 S2 I2 C1 F0}, diff_note={concepts: exponential equation in a logarithmic base | steps: rewrite the powers; substitute; solve; sum}]
If $2x^{\log_43}+3^{\log_4x}=27$, then sum of possible values of $\sqrt x$ is equal to :-
\begin{qoptions}
\opt{A}{0}
\opt{B}{2}
\opt{C}{4}
\opt{D}{6}
\end{qoptions}
\end{bankq}
\begin{bankq}[id=STAR01-P2-MATH-S1-Q14, source={Star Review (Other) - Test 01 (12 Sep 2022)}, paper=2, qno=14, type=single, subtopic={Properties}, answer={B}, answer_src=key, difficulty=5, diff_parts={K1 S2 I1 C2 F0}, diff_note={concepts: number of digits using logarithms | steps: solve the nested logarithms; value of N; logarithm of the product; digits}]
If $\log_5\log_4\log_3\log_2(N)=0$, then number of digits of $N\times2^{19}$ are --\par
\{where $\log_{10}2=0.3010$\}
\begin{qoptions}
\opt{A}{21}
\opt{B}{31}
\opt{C}{20}
\opt{D}{30}
\end{qoptions}
\end{bankq}
\begin{bankq}[id=STAR01-P2-MATH-S1-Q20, source={Star Review (Other) - Test 01 (12 Sep 2022)}, paper=2, qno=20, type=single, subtopic={Equations and inequalities}, answer={A}, answer_src=key, difficulty=5, diff_parts={K1 S2 I2 C1 F0}, diff_note={concepts: exponential and logarithmic equation | steps: simplify the power; quadratic in the logarithm; solve; nature of x}]
$(100)^{\log_{10}(\ell nx)}=\ell nx-(\ell nx)^2+1$, then $x$ is (where $\ell nx=\ell n_ex$)
\begin{qoptions}
\opt{A}{Irrational number}
\opt{B}{Rational number}
\opt{C}{Even integer}
\opt{D}{odd integer}
\end{qoptions}
\end{bankq}
\end{document}
